% Generated draft: requires mathematical and source review.
\documentclass[11pt]{article}
\usepackage[margin=25mm]{geometry}
\usepackage{fontspec}
\setmainfont{Latin Modern Roman}
\usepackage{amsmath,amssymb,mathtools,unicode-math}
\setmathfont{Latin Modern Math}
\usepackage{enumitem,needspace,xurl,hyperref}
\setlength{\parindent}{0pt}
\setlength{\parskip}{5pt}
\setlength{\emergencystretch}{3em}
\newcommand{\N}{\mathbb{N}}
\newcommand{\Z}{\mathbb{Z}}
\newcommand{\Q}{\mathbb{Q}}
\newcommand{\R}{\mathbb{R}}
\newcommand{\C}{\mathbb{C}}
\newcommand{\F}{\mathbb{F}}
\newcommand{\A}{\mathbb{A}}
\newcommand{\PP}{\mathbb{P}}
\newcommand{\E}{\mathbb{E}}
\newcommand{\eps}{\varepsilon}
\newcommand{\dd}{\,\mathrm{d}}
\begin{document}
\section*{UnsolvedMath ID: 10400037}
\textbf{Source problem code:} catalogue entry
\subsection{Definitions and mathematical statement}
Let $I$ be an oriented interval and $S^1$ an oriented circle. Write $\overrightarrow{\mathcal A}(X;\mathbb Q)^{(d)}$ for the degree-$d$ arrow-diagram space used in the source, for $X\in\{S^1,I\}$. In a usual chord presentation an arrow diagram consists of directed chords attached at distinct points of the oriented skeleton, modulo orientation-preserving skeleton equivalence, but the source's full generating/degree conventions are not reproduced.

For every degree $d$, determine all six dimension sequences:
\[
\begin{aligned}
a_{X,d}&=\dim_{\mathbb Q}\overrightarrow{\mathcal A}(X;\mathbb Q)^{(d)},\\
b_{X,d}&=\dim_{\mathbb Q}\left(\overrightarrow{\mathcal A}(X;\mathbb Q)^{(d)}/R_{\mathrm{FI},X,d}\right),\\
c_{X,d}&=\dim_{\mathbb Q}\left(\overrightarrow{\mathcal A}(X;\mathbb Q)^{(d)}/R_{\mathrm{weak\ FI},X,d}\right),
\qquad X\in\{S^1,I\}.
\end{aligned}
\]
Here the two $R$ spaces mean the subspaces generated by the source's oriented FI and weak oriented FI relations in degree $d$, respectively. These are separate quotients.

The extract names the 6T, oriented FI, weak FI, AS, IHX, and STU relations, but their defining pictures in Figure 9 are missing, leaving only signs and equals symbols. Its caption says a thick dashed line represents the sum of its two orientations, and corresponding thin dashed lines in a formula share the same arbitrary orientation. These orientation rules are retained, but they do not reconstruct the lost diagrams. Thus the exact arrow space, its original relations, and both further quotient subspaces must be recovered from Figure 9 before the dimensions refer to fully specified vector spaces. No remembered relation convention or numerical dimension is substituted.
\subsection{Short English statement}
Compute degree-by-degree rational dimensions of the circle and interval arrow-diagram spaces, both before and after the distinct FI and weak-FI quotients. The defining relation pictures are missing, so their exact conventions remain unresolved.
\subsection{Sources}
\begin{itemize}
\item Ohtsuki, \emph{Problems on invariants of knots and 3-manifolds} (2002), Problem 2.15, PDF page 37, and Figure 9 with its orientation conventions, \url{https://msp.org/gtm/2002/04/gtm-2002-04-024s.pdf}; accessed 2026-07-29, PDF-font-structure extraction.
\item Local registry \texttt{data/unsolvedmath/problems.json}, permanent UnsolvedMath ID $10400037$, source code catalogue entry, \emph{Problem 2.15 --- Let I denote an oriented interval}. Dataset attribution: Ulam AI and the UnsolvedMath Contributors.
\end{itemize}
\end{document}
